\section{Results}

\paragraph{The characteristic benchmark.}
% REVIEW: Explain why the effective neighbour count matters. Weights
% spread across nearly the whole fitted roster give each issuer a
% similar peer average. Use that observation to
% motivate fitting an intercept, and keep the common-mean explanation
% conditional.
The joint characteristic kernel has
\ffnum[1]{quarterly-fernandez-joint-effective-neighbours} effective
neighbours per row on average, compared with $n-1$ for a fitted roster of
$n$ issuers under uniform
off-diagonal weights. Its broad peer averages carry little
cross-sectional variation relative to a more concentrated graph. On
the \ffnum[0]{quarterly-primary-issuers} fitted issuers, the
joint-kernel lag maximum-likelihood estimate without an intercept is
$\widehat{\rho}=\ffnum{quarterly-fernandez-joint-rho}$.
Table~\ref{tab:us-benchmark} records the contrast and its
uncertainty. With an intercept, the joint-kernel fit is flagged as
invalid or on the boundary. The no-intercept estimate is consistent
with the spatial term absorbing much of the common mean in this US sample.

\begin{table}[t]
  \centering
  \begin{tabular}{lcc}
    \toprule
    Joint kernel lag ML & $\widehat\rho$ & Refit percentiles \\
    \midrule
    Without intercept & \ffnum{quarterly-fernandez-joint-rho} &
    $[\ffnum{quarterly-fernandez-joint-rho-bootstrap-low},\
    \ffnum{quarterly-fernandez-joint-rho-bootstrap-high}]$
    \\
    With intercept & Unavailable & --- \\
    \bottomrule
  \end{tabular}
  % REVIEW: Explain the 'Unavailable' table cell with the reported
  % invalid-or-boundary diagnostic. Keep the refit percentiles and
  % their conditional interpretation together, and tie incomplete
  % return histories to the data discussion so the two meanings of
  % coverage stay clear.
  \caption{US reimplementation for the December 2015 formation.
    Percentiles are the central 95\% of
    \ffnum[0]{quarterly-fernandez-joint-rho-bootstrap-draws}
    stationary-quarter refits. They describe refit variation
    conditional on the formation graph and the observed-quarter design.
    Issuer availability can change by draw. Statistical coverage of
    this interval
    and complete terminal-outcome coverage of the formation cohort remain
  unverified.}
  \label{tab:us-benchmark}
\end{table}

% REVIEW: Give the two checks distinct subjects. The profile
% calculation checks the point estimate on fixed data, while the
% block refits measure variation when betas and medians are
% re-estimated. Specify which check governs inclusion where the text
% says 'this comparison'.
The point fit and a separate profile-likelihood calculation agree on
the saved original-sample betas, issuer medians and reconstructed
graph. The percentile endpoints recompute betas and medians in common
time blocks and restrict the graph when firms drop out of a refit.
They quantify a different source of variation from the standard error
conditional on the generated beta. We leave estimates for which this
comparison remains unresolved outside the table. These diagnostics
characterize the US reimplementation, where the original sample is unavailable.

\paragraph{Does the encoder add to its inputs?}
On the quarterly panel, the ratio-based TabPFN graph has a higher mean
lag log likelihood than the joint characteristic-history graph when both
use symmetrized ten-neighbour weights. Against peers selected directly
from the financial inputs, the estimated difference changes sign and its
paired percentile interval includes zero (Table~\ref{tab:quarterly-comparisons}). The positive contrast with
the joint graph also falls short of the five-percent threshold after
Holm adjustment. The choice of reference changes what an apparent gain
in spatial fit tells us about the encoder.

\paragraph{A dated view of learned geometry.}
% REVIEW: Bridge from the return fit to the question of which inputs
% align with the principal coordinates. Join the two sample-selection
% sentences with 'because' to explain why an analysis of formation
% inputs can use issuers with incomplete subsequent returns. Name the
% four input histories behind 'original characteristic rung'.
We form principal coordinates from chord distances of the January
2016 TabPFN states on the original characteristic rung. The state and
date match the full analytical formation roster of
\ffnum[0]{quarterly-primary-roster} issuers. This calculation uses
the formation information available to the encoder. It does not apply
the later return-estimation filter. For each ratio we use the latest
available annual slot and restrict the comparison to its observed
% REVIEW: Specify that 'stronger' refers to absolute rank
% correlation, since the coordinate's sign is arbitrary. Give one
% verified valuation comparison and its observed issuer count so
% readers can judge the magnitude and see which observations support
% it.
issuers. The leading coordinate has stronger rank correlation with
market capitalization over assets and market-to-book than with
dividend yield or debt maturity. The same ordering holds for explained
rank variation after regression on the leading three coordinates and
an intercept. This subspace comparison is invariant to sign changes and
orthogonal rotations.

% REVIEW: Lead with what alignment to the valuation inputs reveals
% about this encoding on this date. Keep input overlap beside that
% interpretation and consolidate the repeated pricing and portfolio
% qualifications in the scope appendix.
The valuation ratios are part of the encoder inputs. Their alignment
describes how the model organizes its own inputs on this date.
The historical ratio embedding puts about 45\% of positive inertia
in its leading coordinate. On the 878 issuers with estimated market
betas, its first three coordinates explain about one quarter of one
percent of beta rank variation. The same subspace explains about
93\% of market-cap-to-assets rank variation among the 960 issuers
with that ratio observed.
The evidence leaves open stability across
formations and incremental financial value relative to a
direct-input geometry. A coordinate has an arbitrary orientation and
does not identify a traded factor. The economic pricing restrictions
and model-implied portfolio risk remain open.
